% introduction_taxonomy.tex
\subsection{Taxonomy, Jingle--Jangle Fallacies, and Construct Redundancy}

Embeddings also probe the field's taxonomy: scales with different content share a label (the jingle fallacy); near-interchangeable ones bear different labels (the jangle fallacy). \citet{wulff-mata-2025} co-embedded 4{,}452 IPIP items, 459 scales, and 277 construct labels with a fine-tuned sentence transformer whose similarities predicted empirical internal consistencies ($r = .75$ in sample, $.61$ out of sample), recovered most inventories' factor structures, tracked item and scale intercorrelations, and made misalignment measurable---an automated audit flagged 340 jingle and 164 jangle candidates among IPIP scale pairs, a count that clustering and relabeling more than halved with roughly a quarter as many labels. \citet{wulff-mata-2026} add three applications---mapping measure relations, relabeling taxonomies, and generating new items or constructs vetted in the same space---cautioning that semantic proximity does not settle theoretical identity (subjective and physiological stress may remain conceptually distinct) and that automated consolidation should inform, not replace, expert judgment.

The same geometry detects construct redundancy: \citet{huang-etal-2025} embedded items with a large proprietary model, treated cosine similarities as ``synthetic correlations,'' and applied agglomerative clustering, collecting no responses. Against theory-based benchmarks it separated unrelated constructs---conscientiousness from gratitude, within-cluster similarities far exceeding between-cluster ones---and reproduced a documented redundancy: grit's perseverance-of-effort facet joined conscientiousness, not its own scale, where meta-analytic response data locate the overlap. Output was stable across input configurations and closer to the expected structure than prompted chatbots---reproducible semantic analysis requires an explicit embedding-plus-clustering pipeline, not ad hoc prompting.
